

\chapter{Multicolinéarité}
\label{sec:Multicolinéarité}

\section{Introduction et définition}
Depuis le début du cours, nous avons supposé que la matrice  $\boX$ est une matrice de plein rang, c'est-à-dire qu'il existe une unique inverse à la matrice $\boX'\boX$, et donc qu'il n'existe qu'un seul vecteur $\hbobeta$ qui minimise la somme des carrés résiduels $SC_{\text{res}}$. Toutefois, en pratique certaines colonnes de $\boX$ peuvent être reliées linéairement entre elles. C'est ce qu'on nomme la  multicolinéarité.

%Si certaines colonnes de $X$ sont reliées linéairement entre elles, alors cet inverse n'est plus unique et on ne peut aller plus %loin. Par contre, si certaines colonnes ne sont pas exactement liées de façon linéaire mais sont près de l'être, alors il est %quand même possible de trouver un unique estimateur $\hbobeta$.
La définition mathématique formelle de la multicolinéarité est la suivante.
\begin{defn}\label{def:multi}
Il y a \textbf{multicolinéarité exacte} lorsqu'une variable est la combinaison linéaire d'une ou de plusieurs autres variables, c'est-à-dire  lorsqu'il existe une
 combinaison de constantes $d_1,\ldots,d_p$ (pas toutes \'egales
 \`a z\'ero) telle que
 $$\sum_{j=1}^pd_jX_j=\boz.$$
\end{defn}

\begin{defn}
La \textbf{multicolinéarité} au sens large se produit lorsqu'il existe une
combinaison de constantes $d_1,\ldots,d_p$ (pas toutes \'egales
 \`a z\'ero) telle que
 $$\sum_{j=1}^pd_jX_j\approx \boz.$$
 \end{defn}
Dans le cas de multicolinéarité exacte, la matrice $\boX'\boX$ n'est pas inversible et la régression n'est pas possible.
La multicolinéarité exacte ne représente pas un problème : elle se détecte facilement et il suffit d'enlever la ou les variables en cause pour régler le problème.
La multicolinéarité au sens large est plus fréquente et plus difficile à détecter.

La multicolinéarité au sens large pose plusieurs problèmes :
\begin{enumerate}
  \item Instabilité de $(\boX'\boX)^{-1}$, ce qui entraîne qu'une petite variation en $\boY$ peut se traduire par de très grands changements en $\hbobeta$ et en $\hY$.
  \item On observe des $\hat{\beta}_i$ de signe inattendu ou contre-intuitif.
  \item Les variances des $\hat{\beta}_i$ et des $\hat{Y_i}$ sont très grandes.
  \item Les méthodes de sélection de variables ne concordent pas.
  \item Certains paramètres peuvent se révéler non-significatifs alors que la corrélation des variables correspondantes avec $Y$ est très grande.
\end{enumerate}

Le problème, au fond, est que les variables s'expliquent entre elles. Si on s'imagine l'espace d'une régression à deux variables, c'est comme si les points étaient alignés en «piquets de clôture». Ce phénomène n'est pas facile à détecter dans un espace à haute dimension. Nous présenterons ici quelques outils de détection de la multicolinéarité.

\section{Détection de la multicolinéarité}
\subsection{La matrice de corrélations}
À première vue, on pourrait être tenté de calculer la matrice des corrélations entre les variables exogènes pour détecter la multicolinéarité.

Le coefficient de corrélation échantillonnal entre la variable $x_j$ et la variable $x_k$ est donné par
$$r_{jk}=\dfrac{1}{n-1}\sum_{i=1}^n\frac{(x_{ji}-\overline x_j)(x_{ki}-\overline x_k)}{s_js_k}, \;j\in\{1,\ldots,p\},k\in\{1,\ldots,p\}$$
où $\overline x_j=\sum_{i=1}^n x_{ji}/n$ et $s_j=\sqrt{\sum_{i=1}^n(x_{ji}-\overline x_j)^2/(n-1)}$, et de même pour $\overline x_k$ et $s_k$.

On peut aussi écrire la matrice des coefficients de corrélation échantillonnaux (en excluant la colonne de 1) comme
\begin{eqnarray*}
\boX_c^{*'}\boX_c^{*}=\left(\begin{array}{cccc}
1 & r_{12} &  \cdots & r_{1p} \\
r_{21} & 1 & \cdots & r_{2p} \\
\vdots & \vdots & \ddots  & \vdots \\
r_{p1} & r_{p2} & \cdots  & 1
\end{array}
\right)
\end{eqnarray*}
où $$\boX_c^*=\left(\frac{\bo{x}_1-\overline{\bo{x}}_1}{\sqrt{\sum_i (x_{1i}-\overline x_1)^2} }\;\;\frac{\bo{x}_2-\overline{\bo{x}}_2}{\sqrt{\sum_i (x_{2i}-\overline x_2)^2}} \ldots \frac{\bo{x}_p-\overline{\bo{x}}_p}{\sqrt{\sum_i (x_{pi}-\overline x_p)^2}}  \right)\quad \text{ et } \quad \overline{\bo{x}}_j=(\overline x_j,\ldots,\overline x_j)'.$$



Si deux variables exogènes sont linéairement reliées, leur coefficient de corrélation devrait
être élevé. Cependant deux problèmes majeurs empêchent les coefficients de corrélation
d'être des bons outils de diagnostic:
\begin{itemize}
\item[(i)] il est difficile de dire ce qu'est une grande valeur de la corrélation
\item[(ii)] la multicolinéarité est souvent induite par une dépendance linéaire entre plus de deux variables.
\end{itemize}
Pour voir que le point (ii) constitue un
problème, il est possible de créer des exemples où $p$ variables sont parfaitement linéairement
dépendantes, mais où les coefficients de corrélation de toutes les paires de variables sont
inférieurs à $1/(p-1)$. Ainsi, même si la corrélation entre deux variables est indésirable, cela
ne représente qu'un cas particulier de la multicolinéarité.

\subsection{Le facteur d'inflation de la variance (VIF)}
\label{sec:vif}
Une approche plus sensible consiste à évaluer le degré de dépendance linéaire de chaque
variable exogène sur les autres variables exogènes. Ainsi pour la $j^e$ variable exogène,
on peut mesurer ce niveau de dépendance en effectuant une régression linéaire avec la $j^e$
variable exogène comme variable réponse et les $p-1$ variables exogènes restantes comme variables explicatives. Le coefficient de détermination de cette régression, noté $R^2_j$, mesurera la proportion de la variabilité de la $j^e$ variable exogène qui est expliquée de fa\c{c}on linéaire par les autres variables exogènes. Si $R_j^2$ est grand, on a vraisemblablement un problème.

Cette idée se matérialise dans le facteur d'inflation de la variance (en anglais, VIF pour Variance Inflation Factor). Celui-ci est défini de la façon suivante :
\begin{eqnarray}
  \mbox{VIF}_j &=& \frac{1}{1-R_j^2}
\end{eqnarray}
Si $R_j^2\rightarrow 1$, alors $VIF_j\rightarrow \infty$. Ainsi, une grande valeur de $\mbox{VIF}_j$ signifie une valeur proche de 1 de $R_j^2$, et donc une forte dépendance linéaire des variables exogènes.

Le nom \og facteur d'inflation de la variance \fg  $\;$ vient du fait que la variance de $\hat{\beta}_j$ s'exprime en fonction du VIF comme suit :
\begin{eqnarray*}
% \nonumber to remove numbering (before each equation)
  \var[\hat{\beta}_j] &=& \frac{\sigma^2}{(\boX'\boX)_{jj}} \mbox{VIF}_j
\end{eqnarray*}
On voit donc l'effet de cette multicolinéarité sur la variance des estimateurs des coefficients.

 Plusieurs
auteurs suggèrent un VIF > 10 comme point où l'on doit commencer à considérer la présence
de multicolinéarité aux conséquences importantes. Cependant, les VIF à eux seuls ne sont
pas un diagnostic complet. Parmi les points faibles des VIF on dénote
\begin{itemize}
\item [(i)]l'incapacité de détecter des multicolinéarités impliquant la colonne de 1 de la matrice d'incidence
\item [(ii)] l'incapacité de cerner le nombre de quasi dépendances linéaires présentes dans les données
\item [(iii)] le fait qu'on a jamais réussi à cerner une valeur précise pour le VIF où l'on doit vraiment commencer à s'inquiéter (10 est une valeur \textit{ ad hoc}).
\end{itemize}

\begin{remarque}
Une mesure équivalente au facteur d'inflation de la variance est la tolérance définie comme suit :
$$\mbox{TOL}_j=1-R_j^2=1/\mbox{VIF}_j.$$
La tolérance $\mbox{TOL}_j$ est donc la partie de la variabilité dans la j$^{\text{ème}}$ variable exogène qui n'est pas expliquée par les autres variables exogènes. Ainsi, si on définit un cas problématique comme un cas où $\mbox{VIF}>10$, ceci est équivalent à définir les cas problématiques comme ceux ayant une $\mbox{TOL}< 10 \%$, c'est-à-dire ceux pour qui au moins 90 \% de la variabilité est expliquée par les autres variables exogènes.
\end{remarque}

\subsection{L'indice de conditionnement}

\label{sec:FacteurDeConditionnement}
L'approche la plus efficace afin de détecter les causes de multicolinéarité est basée sur les valeurs propres de la matrice des coefficients de corrélation échantillonnaux $\boX_c^{*'}\boX_c^*$, ou sur la matrice des valeurs réduites, mais non centrées $\boX^*$ (qui permet l'inclusion de la colonne de 1), où 

$$\boX^*=\left(\frac{\bo{x}_1}{\sqrt{\sum_i x_{1i}^2} }\;\;\frac{\bo{x}_2}{\sqrt{\sum_i x_{2i}^2}} \ldots \frac{\bo{x}_p}{\sqrt{\sum_i x_{pi}^2}}  \right).$$

En se servant de la décomposition en valeurs propres de $\boX^{*'}\boX^*$, i.e.
$$\boX^{*'}\boX^*=\bo{P}\bo{\Lambda} \bo{P}'$$
avec $\bo{\Lambda}=\text{diag}(\lambda_1,\ldots,\lambda_{p'})$ et $\bo{P}\bo{P}'= \bo{P}'\bo{P}=\bo{I}_p$, on peut alors montrer que
\begin{equation}
\text{tr}(\boX^{*'}\boX^*)=\text{tr}(\bo{P}\bo{\Lambda} \bo{P}')=\mbox{tr}(\bo{P}'\bo{P}\bo{\Lambda})=\mbox{tr}(\bo{\Lambda})=\sum_{j=1}^{p'}\lambda_j=p'.
\label{eq:traceIC}
\end{equation}
En présence de multicolinéarité, le déterminant de $\boX^{*'}\boX^*$, qui est égal à $\prod\limits_{j=1}^{p'}\lambda_j$, aura une valeur proche de 0.



Ainsi, en présence de multicolinéarité l'équation \ref{eq:traceIC} indique que la somme des valeurs propres doit être égale à $p'$, mais que leur produit sera proche de 0. Puisque les valeurs propres sont toutes positives, cela implique qu'en présence de multicolinéarité, il y aura \`a la fois des valeurs propres de grandes et de faibles valeurs. Le nombre de $\lambda_j$ proche de 0 indique le nombre variables en multicolinéarité. Dans cette situation, le rapport de la plus grande valeur propre sur la plus faible sera donc élevé.

Les \textbf{indices de conditionnement} notés $\eta_j$ pour ${j=1,2,\cdots,p'}$ sont définis comme les racines carrées des rapports entre la plus grande valeur propre de la matrice $\boX^{*'}\boX^*$ et chacune de ses autres valeurs propres, c'est-à-dire
$$\eta_j=\sqrt{\frac{\lambda_{max}}{\lambda_j}}.$$

Plusieurs auteurs semblent s'entendre pour dire qu'une valeur de $\eta_j > 30$ indique une multicolinéarité dont il faut se
soucier.

La plus grande valeur de $\eta_j$ est égale à
$$\kappa=\frac{\lambda_{max}}{\lambda_{min}}.$$

$\kappa$ est parfois appelée \textbf{ facteur de conditionnement}.


%Les valeurs propres $\lambda_1, ..., \lambda_{p'}$ de $(\boX^'\boX)^{-1}$ sont réelles et positives. On peut écrire $(\boX'\boX)=\bo{P}\bo{\Lambda} \bo{P}'$ où $\bo{P}$ est la matrice orthogonale des vecteurs propres et $\bo{Lambda}$ est la matrice diagonale des valeurs propres de $(\boX'\boX)$.
%
%Il y a multicolinéarité lorsque $\lambda_{p'}$, la plus petite valeur propre de $\boX'\boX$, est proche de 0. Rappelons que le déterminant de $\mbox{det}(\boX'\boX)= \prod_{i=1}^{p'} \lambda_{i}=0$ lorsque $\lambda_{p'}=0$. Donc, si $\lambda_{p'}\cong 0$, $(\boX^'\boX)^{-1}$ est instable numériquement.
%
%
% à continuer ici. vérifier la définition du facteur de conditionnement.
%

%Le facteur de conditionnement est un indicateur de la multicolinéarité qui utilise les valeurs propres de la matrice $(\boX^{*'}\boX^*)$. On  le définit de la façon suivante:
%\begin{eqnarray}
%% \nonumber to remove numbering (before each equation)
%  \kappa &=& \sqrt{\frac{\lambda_{\text{max}}}{\lambda_{\text{min}}}}.
%\end{eqnarray}
%Bien que ceci n'ait pas de véritable fondement, on considère généralement un $\kappa>30$ comme étant problématique.

Soit $\bo{p}_j$, le vecteur propre correspondant à la valeur propre $\lambda_j$.
%Les éléments de ces vecteurs peuvent nous aider à identifier le nombre de quasi dépendances linéaires dans les données, et quelles variables exogènes font partie de chaque quasi dépendance. En effet,
Étant donné que
$$\var[\hat\beta_{k-1}]=\sigma^2(\bo{X}^{*'}\bo{X}^*)^{-1}_{kk}=\sigma^2\sum_{j=1}^{p'}p_{kj}^2/\lambda_j,$$
la portion de variance de $\hat\beta_{k-1}$ expliquée par la $j^{e}$ dépendance linéaire peut se mesurer comme suit :
$$v_{kj}=\frac{p_{kj}^2/\lambda_j}{\sum\limits_{l=1}^{p'}p_{kl}^2/\lambda_l}.$$

Pour détecter la multicolinéarité on suit donc la procédure suivante:
\begin{enumerate}
\item Calculer les indices de conditionnement $\eta_j$ et vérifier si certains d'entre eux possèdent une valeur supérieure à 30.
\item Pour chaque $\eta_j$ dont la valeur excède 30, calculer les proportions de la variance $v_{1j},v_{2j},\cdots,v_{p'j}$.
\item Les variables exogènes correspondant aux coefficents avec des valeurs de $v_{kj}$ élevées (en général nettement supérieures à 60 \%) sont multicolinéaires.
\end{enumerate}

\section{Les solutions possibles à la multicolinéarité}
\label{sec:solutionsMulticolinearite}

La solution la plus simple est de retirer une ou plusieurs variables exogènes. Typiquement, on retranchera les variables ayant le plus grand VIF, une à la fois. Dans la pratique, ce retrait des variables doit se faire en collaboration étroite avec le client, car dans le choix des variables, c'est d'abord l'interprétation qui compte. Il est parfois possible  de combiner des variables exogènes redondantes, par exemple pour calculer un score à partir de plusieurs variables exogènes.

Quelques autres solutions sont également disponibles.



\subsection{La régression \og ridge \fg}

La régression ridge, présentée à la section \ref{sec:ridge} peut également être utilisée en cas de multicolinéarité, mais donne des estimateurs biaisés. Ce n'est toutefois pas une option très populaire en pratique.


\subsection{La régression basée sur les composantes principales}


On peut également faire une régression sur les composantes principales de la matrice d'incidence. Les facteurs dans une analyse en composantes principales sont donnés par les colonnes de la matrice $\bo{F}$ suivante :
$\bo{F}=\boX^*\bo{P}$,
où $\bo{P}$ est la matrice des vecteurs propres de $\boX^{*'}\boX^*$. Cette matrice est telle que $\bo{P}'\bo{P}=\bo{P}\bo{P}'=\boI$. En utilisant cette propriété matricielle, nous pouvons écrire le modèle de régression multiple comme
\begin{eqnarray*}
\boY & = & \bo{1}+\boX^*\bobeta + \boeps \\
 & = & \bo{1}+\boX^*\bo{P}\bo{P}'\bobeta+ \boeps\\
& = & \bo{1}+\bo{F}\bo{\gamma}+\epsilon,
\end{eqnarray*}
où $\bo{\gamma}=\bo{P}'\bobeta$ ou encore $\bobeta=\bo{P}\bo{\gamma}$. Le théorème de Gauss-Markov implique que $\hbobeta=\bo{P}\hat{\bo{\gamma}}.$
La multicolinéarité n'est plus présente, mais au prix d'un biais dans les estimateurs. Ce biais est présent si on ne retient dans la modélisation que les facteurs associés aux plus grandes valeurs propres de $\boX^{*'}\boX^*$.

Le principal problème de la régression en composantes principales est son interprétation, qui peut être hasardeuse. %On peut se rattraper en calculant $\hbobeta=\bo{P}\hat{\bo{\gamma}},$ cependant dans ce cas $\esp[\hbobeta]\neq\bobeta$. L'estimation est biaisée.

%Enfin, on peut faire une régression en composantes principales. Pour ce faire, on commence par centrer-réduire la matrice X, on obtient $X^*$. On calcule ensuite la matrice des corrélation $X^{*T}X^*$. On trouve les valeurs propres de cette matrice et on effectue la régression avec le modèle suivant:
%\begin{eqnarray}
%% \nonumber to remove numbering (before each equation)
%  Y &=& Z\alpha + \epsilon
%\end{eqnarray}
%où $Z=X^*P^T$ et $\alpha=P\beta$. On ne retient que le nombre de composantes nécessaires.
%
%Le principal problème de la régression en composantes principales est son interprétation, qui peut être hasardeuse. On peut se rattraper en calculant $\hat{\beta}=P\hat{\alpha}$, mais dans ce cas $E[\hat{\beta}]\neq\beta$; l'estimation est biaisée.

